\documentclass[11pt,a4paper]{article}

\newcommand{\runninghead}{Study 2 — Temperature compensation, legacy era, 2021}
% ---------------------------------------------------------------------------
% reportstyle.tex — shared preamble for the study reports under studies/.
%
% Kept deliberately inside the package set shipped with TeX Live "basic":
% no tcolorbox, no siunitx, no enumitem, no titlesec. Everything below
% compiles with a stock pdflatex installation.
%
% Usage, from a report directory one level below a study folder:
%     % ---------------------------------------------------------------------------
% reportstyle.tex — shared preamble for the study reports under studies/.
%
% Kept deliberately inside the package set shipped with TeX Live "basic":
% no tcolorbox, no siunitx, no enumitem, no titlesec. Everything below
% compiles with a stock pdflatex installation.
%
% Usage, from a report directory one level below a study folder:
%     % ---------------------------------------------------------------------------
% reportstyle.tex — shared preamble for the study reports under studies/.
%
% Kept deliberately inside the package set shipped with TeX Live "basic":
% no tcolorbox, no siunitx, no enumitem, no titlesec. Everything below
% compiles with a stock pdflatex installation.
%
% Usage, from a report directory one level below a study folder:
%     \input{../../_shared/reportstyle.tex}
%     \graphicspath{{../outputs/}}
% ---------------------------------------------------------------------------

\usepackage[T1]{fontenc}
\usepackage[utf8]{inputenc}
\usepackage{textcomp}
\usepackage{lmodern}
\usepackage{microtype}
\usepackage[margin=2.4cm,headheight=15pt]{geometry}
\usepackage{graphicx}
\usepackage{booktabs}
\usepackage{tabularx}
\usepackage{array}
\usepackage{amsmath}
\usepackage{xcolor}
\usepackage[font=small,labelfont=bf,justification=justified]{caption}
\usepackage{float}
\usepackage{fancyhdr}
\usepackage[hidelinks]{hyperref}

% --- Palette ---------------------------------------------------------------
\definecolor{rulegrey}{HTML}{B9C0C7}
\definecolor{fillgrey}{HTML}{F2F4F6}
\definecolor{failred}{HTML}{A5202B}
\definecolor{passgreen}{HTML}{1B6B3A}
\definecolor{accent}{HTML}{1F4E79}

% --- Semantic shorthands ---------------------------------------------------
\newcommand{\degC}{\textdegree C}
\newcommand{\mdegC}{mdeg/\textdegree C}
\newcommand{\fail}{\textcolor{failred}{\textbf{fail}}}
\newcommand{\pass}{\textcolor{passgreen}{\textbf{pass}}}
\newcommand{\worse}{\textcolor{failred}{worse}}
\newcommand{\better}{\textcolor{passgreen}{better}}

% --- Highlighted panel -----------------------------------------------------
% #1 = heading, #2 = body
\newcommand{\panel}[2]{%
  \begin{center}%
  \setlength{\fboxsep}{9pt}%
  \fcolorbox{rulegrey}{fillgrey}{%
    \begin{minipage}{0.93\textwidth}%
      {\bfseries\color{accent}#1}\par\medskip
      #2
    \end{minipage}}%
  \end{center}}

% --- Numeric column: right aligned, fixed width ----------------------------
\newcolumntype{R}[1]{>{\raggedleft\arraybackslash}p{#1}}
\newcolumntype{L}[1]{>{\raggedright\arraybackslash}p{#1}}
\newcolumntype{Y}{>{\raggedright\arraybackslash}X}

% --- Table look ------------------------------------------------------------
\renewcommand{\arraystretch}{1.15}
\setlength{\tabcolsep}{6pt}

% --- Headers ---------------------------------------------------------------
\pagestyle{fancy}
\fancyhf{}
\fancyhead[L]{\small\itshape\runninghead}
\fancyhead[R]{\small\thepage}
\renewcommand{\headrulewidth}{0.4pt}

% --- Section spacing without titlesec --------------------------------------
\makeatletter
\renewcommand\section{\@startsection{section}{1}{\z@}%
  {-3.2ex \@plus -1ex \@minus -.2ex}{2.0ex \@plus.2ex}%
  {\normalfont\Large\bfseries\color{accent}}}
\renewcommand\subsection{\@startsection{subsection}{2}{\z@}%
  {-2.6ex \@plus -1ex \@minus -.2ex}{1.4ex \@plus.2ex}%
  {\normalfont\large\bfseries}}
\makeatother

% --- Figure defaults -------------------------------------------------------
\setkeys{Gin}{width=\linewidth}
\captionsetup[figure]{skip=6pt}

%     \graphicspath{{../outputs/}}
% ---------------------------------------------------------------------------

\usepackage[T1]{fontenc}
\usepackage[utf8]{inputenc}
\usepackage{textcomp}
\usepackage{lmodern}
\usepackage{microtype}
\usepackage[margin=2.4cm,headheight=15pt]{geometry}
\usepackage{graphicx}
\usepackage{booktabs}
\usepackage{tabularx}
\usepackage{array}
\usepackage{amsmath}
\usepackage{xcolor}
\usepackage[font=small,labelfont=bf,justification=justified]{caption}
\usepackage{float}
\usepackage{fancyhdr}
\usepackage[hidelinks]{hyperref}

% --- Palette ---------------------------------------------------------------
\definecolor{rulegrey}{HTML}{B9C0C7}
\definecolor{fillgrey}{HTML}{F2F4F6}
\definecolor{failred}{HTML}{A5202B}
\definecolor{passgreen}{HTML}{1B6B3A}
\definecolor{accent}{HTML}{1F4E79}

% --- Semantic shorthands ---------------------------------------------------
\newcommand{\degC}{\textdegree C}
\newcommand{\mdegC}{mdeg/\textdegree C}
\newcommand{\fail}{\textcolor{failred}{\textbf{fail}}}
\newcommand{\pass}{\textcolor{passgreen}{\textbf{pass}}}
\newcommand{\worse}{\textcolor{failred}{worse}}
\newcommand{\better}{\textcolor{passgreen}{better}}

% --- Highlighted panel -----------------------------------------------------
% #1 = heading, #2 = body
\newcommand{\panel}[2]{%
  \begin{center}%
  \setlength{\fboxsep}{9pt}%
  \fcolorbox{rulegrey}{fillgrey}{%
    \begin{minipage}{0.93\textwidth}%
      {\bfseries\color{accent}#1}\par\medskip
      #2
    \end{minipage}}%
  \end{center}}

% --- Numeric column: right aligned, fixed width ----------------------------
\newcolumntype{R}[1]{>{\raggedleft\arraybackslash}p{#1}}
\newcolumntype{L}[1]{>{\raggedright\arraybackslash}p{#1}}
\newcolumntype{Y}{>{\raggedright\arraybackslash}X}

% --- Table look ------------------------------------------------------------
\renewcommand{\arraystretch}{1.15}
\setlength{\tabcolsep}{6pt}

% --- Headers ---------------------------------------------------------------
\pagestyle{fancy}
\fancyhf{}
\fancyhead[L]{\small\itshape\runninghead}
\fancyhead[R]{\small\thepage}
\renewcommand{\headrulewidth}{0.4pt}

% --- Section spacing without titlesec --------------------------------------
\makeatletter
\renewcommand\section{\@startsection{section}{1}{\z@}%
  {-3.2ex \@plus -1ex \@minus -.2ex}{2.0ex \@plus.2ex}%
  {\normalfont\Large\bfseries\color{accent}}}
\renewcommand\subsection{\@startsection{subsection}{2}{\z@}%
  {-2.6ex \@plus -1ex \@minus -.2ex}{1.4ex \@plus.2ex}%
  {\normalfont\large\bfseries}}
\makeatother

% --- Figure defaults -------------------------------------------------------
\setkeys{Gin}{width=\linewidth}
\captionsetup[figure]{skip=6pt}

%     \graphicspath{{../outputs/}}
% ---------------------------------------------------------------------------

\usepackage[T1]{fontenc}
\usepackage[utf8]{inputenc}
\usepackage{textcomp}
\usepackage{lmodern}
\usepackage{microtype}
\usepackage[margin=2.4cm,headheight=15pt]{geometry}
\usepackage{graphicx}
\usepackage{booktabs}
\usepackage{tabularx}
\usepackage{array}
\usepackage{amsmath}
\usepackage{xcolor}
\usepackage[font=small,labelfont=bf,justification=justified]{caption}
\usepackage{float}
\usepackage{fancyhdr}
\usepackage[hidelinks]{hyperref}

% --- Palette ---------------------------------------------------------------
\definecolor{rulegrey}{HTML}{B9C0C7}
\definecolor{fillgrey}{HTML}{F2F4F6}
\definecolor{failred}{HTML}{A5202B}
\definecolor{passgreen}{HTML}{1B6B3A}
\definecolor{accent}{HTML}{1F4E79}

% --- Semantic shorthands ---------------------------------------------------
\newcommand{\degC}{\textdegree C}
\newcommand{\mdegC}{mdeg/\textdegree C}
\newcommand{\fail}{\textcolor{failred}{\textbf{fail}}}
\newcommand{\pass}{\textcolor{passgreen}{\textbf{pass}}}
\newcommand{\worse}{\textcolor{failred}{worse}}
\newcommand{\better}{\textcolor{passgreen}{better}}

% --- Highlighted panel -----------------------------------------------------
% #1 = heading, #2 = body
\newcommand{\panel}[2]{%
  \begin{center}%
  \setlength{\fboxsep}{9pt}%
  \fcolorbox{rulegrey}{fillgrey}{%
    \begin{minipage}{0.93\textwidth}%
      {\bfseries\color{accent}#1}\par\medskip
      #2
    \end{minipage}}%
  \end{center}}

% --- Numeric column: right aligned, fixed width ----------------------------
\newcolumntype{R}[1]{>{\raggedleft\arraybackslash}p{#1}}
\newcolumntype{L}[1]{>{\raggedright\arraybackslash}p{#1}}
\newcolumntype{Y}{>{\raggedright\arraybackslash}X}

% --- Table look ------------------------------------------------------------
\renewcommand{\arraystretch}{1.15}
\setlength{\tabcolsep}{6pt}

% --- Headers ---------------------------------------------------------------
\pagestyle{fancy}
\fancyhf{}
\fancyhead[L]{\small\itshape\runninghead}
\fancyhead[R]{\small\thepage}
\renewcommand{\headrulewidth}{0.4pt}

% --- Section spacing without titlesec --------------------------------------
\makeatletter
\renewcommand\section{\@startsection{section}{1}{\z@}%
  {-3.2ex \@plus -1ex \@minus -.2ex}{2.0ex \@plus.2ex}%
  {\normalfont\Large\bfseries\color{accent}}}
\renewcommand\subsection{\@startsection{subsection}{2}{\z@}%
  {-2.6ex \@plus -1ex \@minus -.2ex}{1.4ex \@plus.2ex}%
  {\normalfont\large\bfseries}}
\makeatother

% --- Figure defaults -------------------------------------------------------
\setkeys{Gin}{width=\linewidth}
\captionsetup[figure]{skip=6pt}

\graphicspath{{../outputs/}}

\title{\vspace{-1.4cm}\bfseries Is the documented temperature compensation good for
the prediction system?\\[4pt]
\large\mdseries Study 2 --- the legacy instrument era, station 02, calendar year 2021\\
\large\mdseries Medieval walls of Gubbio, static structural health monitoring}
\author{}
\date{Notebook executed 11 August 2026 \quad$\cdot$\quad
Data window 1 January -- 31 December 2021 \quad$\cdot$\quad
Hourly analysis grid}

\begin{document}
\maketitle
\thispagestyle{fancy}

%% ==========================================================================
\section*{Summary}
\addcontentsline{toc}{section}{Summary}

This report asks the same questions as its companion study on the current instrument
era, on different data. The correction under test is the fixed thermal correction that
\texttt{00\_Sensor\_Preprocessing} applies before any other stage sees the signal:
%
\begin{equation}
I_{\text{comp}}(t) \;=\; I(t) \;-\; \bigl(T_{\text{air}}(t) - T_{\text{air}}(t_0)\bigr)
\cdot c \cdot 1000 ,
\qquad c = 0.005 .
\label{eq:comp}
\end{equation}
%
Where the current-era study had four environmental channels but only half an annual
cycle, this one has a \emph{full calendar year} and only two channels. That trade is
the reason to run it: the seasonal response is observable here and was not observable
there.

\panel{Verdict on the correction: it fails all three criteria, on independent data}{%
\begin{tabularx}{\linewidth}{@{}Y R{4.3cm} l@{}}
\toprule
\textbf{Criterion} & \textbf{Measured} & \textbf{Result}\\
\midrule
Thermal dependence is reduced
  & $|r|$: $0.238 \rightarrow 0.738$
  & \fail\\
Coefficient is near the variance-minimising value
  & $5.00$ vs $1.00$ \mdegC
  & \fail\\
Forecast improves under a controlled contrast
  & penalty $+0.7\,\%$ to $+47.4\,\%$
  & \fail\\
\bottomrule
\end{tabularx}%
\medskip

\noindent Three for three, matching the current-era study on a different instrument and
a different set of regressors.}

\panel{The finding only a full year could produce: the structure \emph{leads} the air}{%
Over 2021 the inclination correlates with air temperature at $0.238$ in \textbf{levels}
but $0.866$ in \textbf{first differences}. The reason is a phase difference, and its
direction is the important part. A 30-day mean of the inclination peaks on
\textbf{7 June}; a 30-day mean of air temperature peaks on \textbf{5 August}. The
structure reaches its seasonal maximum \textbf{59 days before} the air does, and a
constrained daily cross-correlation independently puts the lead at 89 days.\par\medskip
\textbf{No thermal lag can represent this.} A causal filter can only delay a driver,
never advance it, and the scan confirms it: every time constant from 1 to 168 hours
makes the annual fit monotonically \emph{worse}, from $R^2 = 0.056$ at $\tau = 0$ down
to $0.000$. The annual movement of this wall is not a delayed response to air
temperature at all.\par\medskip
It is, however, what the physical hypothesis predicts. If the forcing is solar gain on
an exposed vertical face, its annual maximum arrives well before the annual maximum of
air temperature, because irradiance on a vertical surface peaks when the sun is lower
than at the solstice. Air temperature is a good proxy for the forcing from hour to hour
--- at the diurnal scale it explains $91.5\,\%$ of the variance, instantaneously --- and
a poor one from month to month.}

\noindent\textbf{Recommendation.} Remove the fixed compensation from Notebook 00. Beyond
that, the two timescales need different treatment: the daily response can be carried by
instantaneous air temperature, while the seasonal component must be left to the
decomposition's own annual term or driven by a radiation-like variable, because no
coefficient and no lag on air temperature can reproduce a lead.

\medskip
{\small\tableofcontents}

\clearpage
%% ==========================================================================
\section{What differs from the current-era study}

\begin{table}[H]
\centering
\caption{The two studies, side by side.}
\begin{tabularx}{\linewidth}{@{}L{4.2cm} Y Y@{}}
\toprule
& \textbf{Current era (Study 1)} & \textbf{Legacy era (this study)}\\
\midrule
Period & from 2025-02-21 & before 2025-02-21\\
Environmental channels & \texttt{tair}, \texttt{rh}, \texttt{sr}, \texttt{twall} & \texttt{tair}, \texttt{rh} only\\
Usable span & 197 days in two blocks & a full calendar year\\
Seasonal response & not observable & observable\\
Gap structure & two multi-month outages & 84 gaps, none over 24 hours\\
Instrument & installed 2025-02-21 & legacy block \texttt{b2}\\
\bottomrule
\end{tabularx}
\end{table}

The legacy record shares no baseline with the current one: the 2025 installation was a
physical reinstallation, so the absolute levels are unrelated. That is not a limitation
here, because every quantity this report relies on --- correlations, slopes, variances,
forecast errors --- is invariant to the arbitrary offset.

\medskip
\noindent\textbf{Sampling.} As in the companion study, the record is aggregated from its
native 20-minute rate to one hour immediately after parsing, matching the hourly rate of
every external proxy the pipeline aligns against. An hour is kept only if at least two
of its three raw samples survived cleaning.

\medskip
\noindent\textbf{Shared tests.} The acceptance tests, coefficient sweep, slope fitting,
lag and inertia analysis and forecasting experiment are \emph{imported} from the sibling
study rather than reimplemented: \texttt{lc\_lib.py} adds only the legacy loader, the
year-completeness scan, the decomposition and the imputation. Any difference in results
between the two reports is therefore a difference in the data, not in the method.

%% ==========================================================================
\section{Choosing the year}

Whole-day file counts mislead here, because a day can exist as a file while station
02's own block is written as zeros throughout it. The scan measures what the study
actually needs: hourly intervals in which the inclinometer, air temperature and
humidity are \emph{all} present.

Only three years required measuring. From the archive audit, 2018 covers 159 days
(acquisition began in July), 2022, 2023 and 2024 miss 212, 171 and 82 whole days
respectively, and the legacy record ends 51 days into 2025.

\begin{table}[H]
\centering
\caption{Usable completeness by year, station 02, on the hourly grid.}
\begin{tabular}{@{}l R{2.0cm} R{2.2cm} R{2.6cm} R{3.4cm} l@{}}
\toprule
\textbf{Year} & \textbf{Hours} & \textbf{Complete} & \textbf{Complete [\%]} &
\textbf{Longest run [days]} & \textbf{Verdict}\\
\midrule
\textbf{2021} & 8\,737 & 8\,516 & $\mathbf{97.47}$ & 66.6 & \textbf{selected}\\
2019 & 8\,737 & 7\,505 & $85.90$ & 34.0 & \\
2020 & 8\,761 & 7\,270 & $82.98$ & 43.8 & \\
\bottomrule
\end{tabular}
\end{table}

\begin{figure}[H]
\includegraphics{S2_02_st02_year_completeness.png}
\caption{Left, per-channel availability by year --- the three channels fail and recover
together, so the losses are acquisition outages rather than individual sensor faults.
Right, usable completeness, the quantity that decides the selection.}
\label{fig:years}
\end{figure}

The equal bar heights in the left panel are informative in themselves: the three
channels are available in exactly the same proportion in every year, which is the
signature of whole-station dropouts rather than channel-specific failure. Nothing needs
to be modelled channel by channel.

2021 is not merely the most complete year, it is the best-shaped one, and that matters
more than the headline percentage.

\begin{table}[H]
\centering
\caption{Gap profile of the inclinometer channel in 2021, on the hourly grid. The other
two channels are identical.}
\begin{tabular}{@{}l R{2.4cm} R{2.4cm} R{2.4cm}@{}}
\toprule
\textbf{Gap length} & \textbf{Gaps} & \textbf{Hours lost} & \textbf{Share of year}\\
\midrule
Single hour & 44 & 44 & $0.5\,\%$\\
Up to 6 h & 33 & 108 & $1.2\,\%$\\
6--24 h & 7 & 69 & $0.8\,\%$\\
\midrule
\textbf{Total} & \textbf{84} & \textbf{221} & $\mathbf{2.5\,\%}$\\
\bottomrule
\end{tabular}
\end{table}

\textbf{No gap in 2021 exceeds 24 hours}, and the longest unbroken run is 66.6 days.
That is what makes the NeuralProphet imputation in Section~4 genuine interpolation
between known values rather than extrapolation from a seasonal model.

%% ==========================================================================
\section{The year as recorded}

\begin{table}[H]
\centering
\caption{Channel coverage over 2021 (8\,737 hourly intervals).}
\begin{tabular}{@{}l R{2.0cm} R{2.0cm} R{1.8cm} R{1.8cm} R{1.8cm} R{1.8cm}@{}}
\toprule
\textbf{Channel} & \textbf{Present} & \textbf{Available} & \textbf{Min} & \textbf{Max}
& \textbf{Mean} & \textbf{Std}\\
\midrule
\texttt{batt} & 8\,516 & 97.5\,\% & 3.400 & 4.690 & 3.891 & 0.297\\
\texttt{tair} & 8\,516 & 97.5\,\% & $-4.63$ & 38.48 & 13.710 & 8.818\\
\texttt{rh}   & 8\,516 & 97.5\,\% & 19.65 & 99.94 & 71.908 & 22.111\\
\texttt{inc}  & 8\,516 & 97.5\,\% & $-338.67$ & $-205.42$ & $-272.46$ & 33.891\\
\bottomrule
\end{tabular}
\end{table}

The inclination standard deviation over the full year, $33.9$ mdeg, is more than twice
the $14.6$ mdeg measured over the current era's Window A. The difference is not
instrumental --- it is the annual cycle, which Window A only samples in part. The spike
screen removed nothing from this year: with a robust $\sigma$ of $43.3$ mdeg the
threshold sits at $346$ mdeg, and no hourly value approaches it.

\begin{figure}[H]
\includegraphics{S2_05_st02_overview.png}
\caption{The 2021 record: inclination, air temperature, relative humidity and battery
voltage on a common time axis.}
\label{fig:ov}
\end{figure}

Figure~\ref{fig:ov} shows the shape the rest of the report analyses, and the central
fact is already visible. The inclination rises through spring, reaches its highest
values in late May and June, and then turns down --- while air temperature is still
climbing towards its August maximum. The two series do not peak together. Section~5
measures the difference.

%% ==========================================================================
\section{Decomposition and imputation}

\subsection{NeuralProphet decomposition of every series}

Each series is fitted \emph{without} autoregression, so that every component is a
function of time alone. That restriction is deliberate: it keeps the components
interpretable, and it is what allows the same fit to be used for imputation, since a
model with autoregressive terms cannot produce a value inside a gap.

\begin{figure}[H]
\includegraphics[width=0.78\linewidth]{S2_06_st02_decomposition_inc.png}
\caption{Inclination, decomposed. From the top: observed and fitted, trend, annual
term, weekly term, daily term.}
\label{fig:decinc}
\end{figure}

\begin{figure}[H]
\includegraphics[width=0.78\linewidth]{S2_06_st02_decomposition_tair.png}
\caption{Air temperature, decomposed on the same basis.}
\label{fig:dectair}
\end{figure}

\begin{figure}[H]
\includegraphics[width=0.78\linewidth]{S2_06_st02_decomposition_rh.png}
\caption{Relative humidity, decomposed on the same basis.}
\label{fig:decrh}
\end{figure}

\begin{table}[H]
\centering
\caption{Decomposition summary. Component standard deviations are in the native units
of each series.}
\small
\begin{tabular}{@{}l R{1.9cm} R{1.9cm} R{1.8cm} R{1.5cm} R{1.5cm} R{1.5cm} R{1.5cm}@{}}
\toprule
\textbf{Series} & \textbf{Var.\ obs.} & \textbf{Var.\ resid.} & \textbf{Explained} &
\textbf{Trend} & \textbf{Yearly} & \textbf{Weekly} & \textbf{Daily}\\
\midrule
\texttt{inc}  & 1148.583 & 87.735  & 92.4\,\% & 14.035 & \textbf{36.131} & 0.487 & 6.438\\
\texttt{tair} & 77.750   & 10.399  & 86.6\,\% & 2.551  & 8.398 & 0.155 & 2.554\\
\texttt{rh}   & 488.902  & 208.520 & 57.3\,\% & 6.318  & 20.124 & 1.358 & 7.756\\
\bottomrule
\end{tabular}
\end{table}

Three things follow.

\begin{description}
\item[The inclination is an annual signal first and a daily signal second.] Its yearly
term has a standard deviation of $36.1$ mdeg against $6.4$ mdeg for the daily term ---
a factor of five and a half. Trend and seasonality together explain $92.4\,\%$ of its
variance. For a structure whose monitoring question is slow drift, that ratio is the
most important number in the decomposition: the annual environmental cycle is the
dominant thing any anomaly must be detected against.
\item[The weekly term is negligible.] $0.487$ mdeg for the inclination and
$0.155$\,\degC\ for air temperature, against daily terms thirteen to sixteen times
larger. A structure with no human loading cycle and a rural weather station should have
no weekly signature, and they do not. This is a free check that the decomposition is
not manufacturing components.
\item[Humidity is the least deterministic channel.] Only $57.3\,\%$ of its variance is
explained by time alone, against $86.6\,\%$ for temperature. Humidity carries more
weather-event content and less calendar content, which is worth remembering when it is
used as a regressor.
\end{description}

\medskip
\noindent\textbf{A caution on the trend panel, and a consequence.} NeuralProphet's
split between a piecewise-linear trend and a Fourier annual term is not unique, and the
two absorb each other: refitting the same series on a different sampling grid moved the
inclination's trend standard deviation between $10$ and $47$ mdeg across runs of this
study, with the annual term moving to compensate. The trend panel must therefore not be
read as structural drift. It also means the fitted \texttt{season\_yearly} component is
the wrong instrument for measuring seasonal \emph{phase}, which is why Section~5 uses a
plain 30-day rolling mean for that purpose instead.

\subsection{Imputation of the regressors}

Only the regressors are filled, and only with their own trend-plus-seasonal fit --- 221
hours each, $2.5\,\%$ of the year. The inclinometer keeps its gaps throughout the
descriptive analysis: imputing the target and then scoring a forecast against the
imputed values would measure the imputation model rather than the forecast.

\begin{figure}[H]
\includegraphics{S2_08_st02_imputation_tair.png}
\caption{Air temperature: the full year with imputed values overlaid (top), and the
largest single gap with three days of context either side (bottom).}
\label{fig:imptair}
\end{figure}

\begin{figure}[H]
\includegraphics{S2_08_st02_imputation_rh.png}
\caption{Relative humidity, filled on the same basis.}
\label{fig:imprh}
\end{figure}

The detail panel is the honest view of what this imputation can and cannot do. Across
the largest gap the fill reproduces the diurnal shape and joins smoothly to the observed
data at both ends, which is what a trend-plus-seasonal model should achieve. What it
cannot reproduce is that day's particular weather: the filled segment is a smooth
climatological day, while the surrounding observed days carry sharp weather-driven
excursions of several degrees. Filled values are therefore systematically less variable
than observed ones.

%% ==========================================================================
\section{Is the response instantaneous?}
\label{sec:lag}

This is the section the full year makes possible, and its answer differs by timescale.

\subsection{Correlation structure: levels against differences}

\begin{table}[H]
\centering
\caption{Pearson correlations over 2021, on the hourly grid.}
\begin{tabular}{@{}l R{3.2cm} R{3.6cm}@{}}
\toprule
\textbf{Pair} & \textbf{Levels} & \textbf{First differences}\\
\midrule
\texttt{inc} -- \texttt{tair} & $+0.238$ & $+0.866$\\
\texttt{inc} -- \texttt{rh}   & $-0.405$ & $-0.696$\\
\texttt{tair} -- \texttt{rh}  & $-0.692$ & $-0.706$\\
\bottomrule
\end{tabular}
\end{table}

\begin{figure}[H]
\includegraphics{S2_11_st02_corr_heatmaps.png}
\caption{Correlation structure over 2021.}
\label{fig:heat}
\end{figure}

Ordinarily a correlation in levels between two trending series is \emph{inflated} by the
shared trend and falls under differencing. The inclination--temperature pair does the
opposite, and by a wide margin: $0.238$ in levels against $0.866$ in differences. A
relationship that is weak in levels and strong in differences is the signature of two
series that move together at short timescales while their slow components are out of
step.

\subsection{Lag structure at the daily scale}

\begin{figure}[H]
\includegraphics{S2_13_st02_cross_correlation.png}
\caption{Cross-correlation of the inclination against each driver over $\pm 120$ hours,
on levels (left) and on first differences (right).}
\label{fig:ccf}
\end{figure}

\begin{table}[H]
\centering
\caption{Peak of each cross-correlation function at the daily scale.}
\begin{tabular}{@{}l R{2.4cm} R{2.2cm} R{2.6cm} R{2.2cm}@{}}
\toprule
& \multicolumn{2}{c}{\textbf{Levels}} & \multicolumn{2}{c}{\textbf{First differences}}\\
\cmidrule(lr){2-3}\cmidrule(lr){4-5}
\textbf{Driver} & \textbf{Peak $r$} & \textbf{Lag} & \textbf{Peak $r$} & \textbf{Lag}\\
\midrule
\texttt{tair} & $+0.238$ & $0$ h & $+0.866$ & $0$ h\\
\texttt{rh}   & $-0.405$ & $0$ h & $-0.696$ & $0$ h\\
\bottomrule
\end{tabular}
\end{table}

Both drivers peak exactly at zero lag on both views. At the daily scale the response is
contemporaneous, and the instantaneous form of Equation~\eqref{eq:comp} is not the
problem --- the same conclusion the current-era study reached from richer data.

\subsection{Thermal inertia}

\begin{figure}[H]
\includegraphics{S2_36_st02_inertia_scan.png}
\caption{Explained variance and fitted slope against thermal time constant, full
signal, 2021.}
\label{fig:tau}
\end{figure}

\begin{figure}[H]
\includegraphics{S2_37_st02_inertia_scan_band.png}
\caption{The same scan after removing a centred weekly mean, which restricts the
comparison to the diurnal band.}
\label{fig:tauband}
\end{figure}

\begin{table}[H]
\centering
\caption{Best time constant per driver. Every scan returns zero, and every non-zero
time constant makes the fit worse.}
\small
\begin{tabular}{@{}l R{1.7cm} R{1.9cm} R{1.9cm} R{1.7cm} R{1.9cm} R{1.9cm}@{}}
\toprule
& \multicolumn{3}{c}{\textbf{Full signal}} & \multicolumn{3}{c}{\textbf{Diurnal band}}\\
\cmidrule(lr){2-4}\cmidrule(lr){5-7}
\textbf{Driver} & \textbf{Best $\tau$} & \textbf{$R^2$, $\tau{=}0$} & \textbf{$R^2$, best}
& \textbf{Best $\tau$} & \textbf{$R^2$, $\tau{=}0$} & \textbf{$R^2$, best}\\
\midrule
\texttt{tair} & 0 h & $0.056$ & $0.056$ & 0 h & $\mathbf{0.915}$ & $\mathbf{0.915}$\\
\texttt{rh}   & 0 h & $0.164$ & $0.164$ & 0 h & $0.378$ & $0.378$\\
\bottomrule
\end{tabular}
\end{table}

The contrast between the two halves of this table is the study's central measurement.
At the \textbf{diurnal} scale, instantaneous air temperature explains $91.5\,\%$ of the
variance of the inclination --- a very tight, undelayed relationship, matching the
$93.2\,\%$ measured in the current era on a different instrument. At the \textbf{annual}
scale it explains $5.6\,\%$. Adding a thermal time constant does not repair the annual
fit; it degrades it monotonically, from $0.056$ at $\tau = 0$ to $0.000$ at
$\tau = 168$ hours.

\subsection{The seasonal scale: the structure leads the air}

A lag scan reaching five days cannot address an annual relationship, and the obvious
extension has a trap in it: the cross-correlation of two near-annual cycles is itself
near-annual, so the largest \emph{absolute} correlation sits about half a year away
where the two are simply in anti-phase. That is a property of periodicity, not a lag.
Two estimators are therefore used, both immune to it.

\begin{table}[H]
\centering
\caption{Seasonal phase, measured two ways. The rolling-mean estimator assumes nothing
about the shape of the cycle; the cross-correlation is restricted to lags carrying the
sign of the contemporaneous correlation.}
\begin{tabular}{@{}l R{2.8cm} R{2.6cm} R{3.6cm}@{}}
\toprule
\textbf{Series} & \textbf{Peak of 30-day mean} & \textbf{Seasonal range} & \textbf{Lead over \texttt{tair}}\\
\midrule
\texttt{inc}  & 7 June 2021 & $90.0$ mdeg & $\mathbf{59}$ \textbf{days}\\
\texttt{tair} & 5 August 2021 & $24.8$\,\degC & ---\\
\texttt{rh}   & 11 January 2021 & $47.6$ \% & 206 days\\
\midrule
\multicolumn{4}{@{}l}{\itshape Constrained cross-correlation of daily means, $\pm 150$ days}\\
\texttt{inc} vs \texttt{tair} & peak $r = +0.798$ & at $-89$ days & $r = +0.164$ at lag 0\\
\bottomrule
\end{tabular}
\end{table}

\begin{figure}[H]
\includegraphics{S2_40_st02_cross_correlation_seasonal.png}
\caption{Cross-correlation of daily means over $\pm 150$ days. The near-annual
oscillation of the curve itself is the periodicity artefact described above; the
physically meaningful feature is that the positive peak sits at a negative lag.}
\label{fig:ccfseason}
\end{figure}

Both estimators agree on the sign and on the order of magnitude: \textbf{the inclination
reaches its seasonal maximum roughly two to three months before air temperature does}.
Three consequences follow, and they are the reason this study exists.

\begin{description}
\item[No thermal lag can produce this.] A causal filter delays; it never advances. That
is why the inertia scan finds $\tau = 0$ and why every larger value is worse: the model
family being searched cannot represent a lead, so the best it can do is not to move.
\item[The annual movement is not a response to air temperature.] Whatever drives it
arrives before the air does. Solar gain on an exposed face is the natural candidate,
because irradiance on a vertical surface peaks well before the solstice, when the sun is
lower and strikes the wall more squarely --- and the geometry of an embankment, with one
face exposed and one restrained, is exactly the configuration that converts that gain
into a tilt. This study cannot confirm it: the legacy era has no radiometer. It can only
establish that air temperature is the wrong driver at this timescale, and say what to
measure next.
\item[One instantaneous coefficient cannot serve both timescales.] The daily response is
tight, contemporaneous and worth $R^2 = 0.92$; the annual response is loose and
phase-shifted. Equation~\eqref{eq:comp} applies a single constant to both.
\end{description}

%% ==========================================================================
\section{How large is the thermal response over a full year?}

\begin{table}[H]
\centering
\caption{Univariate regression of the inclination on each driver over 2021, with
Newey--West standard errors. The documented coefficient implies $5.0$ \mdegC.}
\begin{tabular}{@{}l R{2.4cm} R{3.6cm} R{1.8cm} R{1.6cm} R{1.8cm}@{}}
\toprule
\textbf{Driver} & \textbf{Slope} & \textbf{95\,\% CI} & \textbf{SE} & \textbf{$R^2$} & \textbf{$n$}\\
\midrule
\multicolumn{6}{@{}l}{\itshape Levels}\\
\texttt{tair} & $0.9139$ & $0.6287$ -- $1.1991$ & $0.1455$ & $0.057$ & 8\,516\\
\texttt{rh}   & $-0.6201$ & $-0.7355$ -- $-0.5046$ & $0.0589$ & $0.164$ & 8\,516\\[2pt]
\multicolumn{6}{@{}l}{\itshape First differences}\\
\texttt{tair} & $2.2603$ & $2.1804$ -- $2.3402$ & $0.0408$ & $0.750$ & 8\,432\\
\texttt{rh}   & $-0.4729$ & $-0.5003$ -- $-0.4455$ & $0.0140$ & $0.484$ & 8\,432\\
\bottomrule
\end{tabular}
\end{table}

The two timescales give $0.91$ and $2.26$ \mdegC\ --- a factor of $2.5$ between them,
and neither near $5.0$. The $R^2$ values tell the same story as Section~\ref{sec:lag}:
instantaneous air temperature is a poor explanation of \emph{where} the inclination sits
over a year ($0.057$) and an excellent explanation of \emph{how it is moving} from hour
to hour ($0.750$).

\begin{figure}[H]
\includegraphics{S2_16_st02_thermal_scatter.png}
\caption{Inclination against air temperature over 2021, with the fitted slope and the
slope implied by the documented coefficient.}
\label{fig:scatter}
\end{figure}

Figure~\ref{fig:scatter} is what a phase-shifted annual cycle looks like in a scatter
plot: not a line with scatter about it but a broad loop. The same air temperature
corresponds to inclinations tens of millidegrees apart depending on whether the year is
warming or cooling at that moment. The relationship is genuinely multi-valued in levels,
so no instantaneous scalar correction, at any coefficient, can collapse it.

%% ==========================================================================
\section{Acceptance tests}

\begin{table}[H]
\centering
\caption{Acceptance tests over 2021. The target for every row is a \emph{lower} value
after correction. Six of seven fail.}
\begin{tabular}{@{}l R{2.4cm} R{2.8cm} R{2.2cm} l@{}}
\toprule
\textbf{Test} & \textbf{Raw} & \textbf{Compensated} & \textbf{Change} & \textbf{Verdict}\\
\midrule
Variance [mdeg$^2$]                    & $1148.58$ & $2381.76$ & $+107.4\,\%$ & \worse\\
Standard deviation [mdeg]              & $33.891$ & $48.803$ & $+44.0\,\%$ & \worse\\
Mean diurnal amplitude [mdeg]          & $21.039$ & $23.920$ & $+13.7\,\%$ & \worse\\
$|r|$ with \texttt{tair}, levels       & $0.2378$ & $0.7383$ & $+210.5\,\%$ & \worse\\
$|r|$ with \texttt{tair}, differenced  & $0.8657$ & $0.9026$ & $+4.3\,\%$ & \worse\\
$|r|$ with \texttt{rh}, levels         & $0.4045$ & $0.3419$ & $-15.5\,\%$ & \better\\
$|r|$ with \texttt{rh}, differenced    & $0.6960$ & $0.6972$ & $+0.2\,\%$ & \worse\\
\bottomrule
\end{tabular}
\end{table}

\begin{table}[H]
\centering
\caption{Signed correlations before and after correction. Both drivers reverse sign.}
\begin{tabular}{@{}l R{2.6cm} R{3.2cm} c@{}}
\toprule
\textbf{Driver} & \textbf{Raw} & \textbf{Compensated} & \textbf{Sign reversed}\\
\midrule
\texttt{tair} & $+0.2378$ & $-0.7383$ & yes\\
\texttt{rh}   & $-0.4045$ & $+0.3419$ & yes\\
\bottomrule
\end{tabular}
\end{table}

The single apparent improvement is again an artefact of passing through zero on the way
to a large correlation of the opposite sign. The correction more than doubles the
variance of the series it is supposed to clean.

\begin{figure}[H]
\includegraphics{S2_19_st02_compensation_effect.png}
\caption{Effect of the documented correction over 2021. Top left, the full year before
and after; top right, the first week in detail; bottom left, the thermal dependence;
bottom right, the mean diurnal cycle.}
\label{fig:effect}
\end{figure}

\begin{description}
\item[Full year.] The raw series spans roughly $110$ mdeg peak to trough. The
compensated series spans more than $200$, plunging to about $-170$ mdeg at the end of
summer before recovering. The correction did not remove a seasonal excursion; it roughly
doubled one and inverted its phase.
\item[First week.] In January the raw and compensated traces are already near mirror
images at daily scale.
\item[Thermal dependence.] The raw cloud is the broad loop of Figure~\ref{fig:scatter}.
The compensated cloud is a tighter, steeply descending band --- a stronger dependence on
temperature than before, not a weaker one.
\item[Mean diurnal cycle.] The raw profile rises from about $-10$ mdeg overnight to
$+7$ mdeg at midday. The compensated profile is inverted and larger.
\end{description}

%% ==========================================================================
\section{Where the documented coefficient sits}

\begin{table}[H]
\centering
\caption{Selected points from the 2021 coefficient sweep.}
\begin{tabular}{@{}l R{3.0cm} R{3.4cm} R{3.0cm}@{}}
\toprule
& \textbf{Coefficient [\mdegC]} & \textbf{Residual variance [mdeg$^2$]} & \textbf{$|r|$ with \texttt{tair}}\\
\midrule
                        & $-1.0$ & $1368.45$ & $0.456$\\
No compensation         & $0.0$  & $1148.58$ & $0.238$\\
                        & $0.8$  & $1084.65$ & $0.031$\\
\textbf{Variance-minimising} & $\mathbf{1.0}$ & $\mathbf{1084.22}$ & $\mathbf{0.023}$\\
                        & $1.2$  & $1090.00$ & $0.076$\\
                        & $2.0$  & $1175.35$ & $0.279$\\
                        & $3.0$  & $1421.99$ & $0.488$\\
\textbf{Documented}     & $\mathbf{5.0}$ & $\mathbf{2381.76}$ & $\mathbf{0.738}$\\
                        & $7.0$  & $3963.53$ & $0.852$\\
\bottomrule
\end{tabular}
\end{table}

\begin{figure}[H]
\includegraphics{S2_21_st02_coefficient_sweep.png}
\caption{Residual variance and residual thermal correlation against the compensation
coefficient over 2021, with the documented value marked.}
\label{fig:sweep}
\end{figure}

Two readings, the second specific to the full-year setting.

First, the documented coefficient is five times the variance-minimising one and sits far
up the ascending branch. Applying it more than doubles the residual variance relative to
leaving the signal alone --- $2382$ against $1149$ mdeg$^2$.

Second, \textbf{even the optimal coefficient buys very little here}. The minimum,
$1084$ mdeg$^2$, is under six per cent below the do-nothing value. In the current era
the same sweep cut residual variance by seventy per cent. The reason is
Section~\ref{sec:lag}: over a full year the dominant thermal component is out of phase
with instantaneous temperature, and an instantaneous linear term cannot remove what is
not aligned with it. The residual thermal correlation still falls almost to zero at the
optimum ($0.023$), so the correction can null the \emph{correlation} while barely
touching the \emph{variance} --- a useful reminder that those are different targets, and
that a compensation tuned on one may be worth almost nothing on the other.

%% ==========================================================================
\section{The prediction experiment}

The configurations, the raw-scale scoring rule and the ridge reference model are those
of the companion study, with \texttt{tair} and \texttt{rh} as the only available
regressors. C against D is the controlled contrast. Predictions of a compensated target
have the correction added back before the error is computed, so all configurations are
scored on the same physical quantity.

\medskip
\noindent\textbf{E and F degenerate here, and that is a result.} Configuration E shifts
each driver to its measured best lag and F filters each through its measured best time
constant. In this era both measurements returned zero for both drivers, so E and F are
\emph{identical to A} by construction and score identically. Reporting them is the point:
the physics-motivated configurations were given every chance to differ and the data
declined to make them differ, at the timescale a forecasting model operates on.

\begin{table}[H]
\centering
\caption{2021, one-hour-ahead ridge forecast, 5\,436 training and 2\,330 test samples.}
\small
\begin{tabular}{@{}l R{2.4cm} R{2.4cm} R{2.4cm}@{}}
\toprule
\textbf{Configuration} & \textbf{MAE [mdeg]} & \textbf{RMSE [mdeg]} & \textbf{vs best}\\
\midrule
A, E, F $\cdot$ raw + regressors      & $\mathbf{0.9878}$ & $1.5088$ & \textbf{best}\\
D $\cdot$ raw, no regressors          & $1.0480$ & $1.5494$ & $+6.1\,\%$\\
C $\cdot$ compensated, no regressors  & $1.5443$ & $2.2041$ & $+56.3\,\%$\\
B $\cdot$ compensated + regressors    & $2.0760$ & $2.6780$ & $+110.2\,\%$\\
\bottomrule
\end{tabular}
\end{table}

\begin{figure}[H]
\includegraphics{S2_23_st02_experiment_h1.png}
\caption{One-hour-ahead forecast, ridge model, 2021.}
\label{fig:h1}
\end{figure}

\begin{table}[H]
\centering
\caption{2021, 24-hour-ahead ridge forecast, 5\,350 training and 2\,294 test samples.}
\small
\begin{tabular}{@{}l R{2.4cm} R{2.4cm} R{2.4cm}@{}}
\toprule
\textbf{Configuration} & \textbf{MAE [mdeg]} & \textbf{RMSE [mdeg]} & \textbf{vs best}\\
\midrule
A, E, F $\cdot$ raw + regressors      & $\mathbf{3.9242}$ & $5.1083$ & \textbf{best}\\
D $\cdot$ raw, no regressors          & $4.0404$ & $5.2348$ & $+3.0\,\%$\\
C $\cdot$ compensated, no regressors  & $4.0692$ & $5.2063$ & $+3.7\,\%$\\
B $\cdot$ compensated + regressors    & $7.9082$ & $8.8802$ & $+101.5\,\%$\\
\bottomrule
\end{tabular}
\end{table}

\begin{figure}[H]
\includegraphics{S2_25_st02_experiment_h24.png}
\caption{Twenty-four-hour-ahead forecast, ridge model, 2021.}
\label{fig:h24}
\end{figure}

Unlike the current-era study, supplying the drivers helps here --- A beats D at both
horizons. That is what a full annual cycle buys: with a whole year in the training split
the model is not forced to extrapolate outside its training temperature range, and the
confound that made the current-era A-versus-D comparison uninformative is absent.

\begin{table}[H]
\centering
\caption{The controlled contrast, C against D.}
\begin{tabular}{@{}l R{2.6cm} R{3.0cm} R{2.6cm}@{}}
\toprule
\textbf{Horizon} & \textbf{MAE raw (D)} & \textbf{MAE compensated (C)} & \textbf{Penalty}\\
\midrule
1 h  & $1.0480$ & $1.5443$ & $+47.4\,\%$\\
24 h & $4.0404$ & $4.0692$ & $+0.7\,\%$\\
\bottomrule
\end{tabular}
\end{table}

The penalty is large at one hour and almost nil at twenty-four. The reason is the same
one that made the coefficient sweep so flat: at the daily horizon the error is dominated
by the part of the signal that an instantaneous thermal term cannot describe either way,
so compensating the target changes little. It still changes it for the worse, and
configuration B --- compensated \emph{and} given the regressors --- remains catastrophic
at $+101.5\,\%$, the double-correction signature seen in the current era.

\subsection{Confirmation under NeuralProphet}

\begin{table}[H]
\centering
\caption{2021, one hour ahead under NeuralProphet, 2\,600 scored points per
configuration, fixed random seed.}
\small
\begin{tabular}{@{}l R{2.4cm} R{2.4cm} R{2.4cm}@{}}
\toprule
\textbf{Configuration} & \textbf{MAE [mdeg]} & \textbf{RMSE [mdeg]} & \textbf{vs best}\\
\midrule
A, E, F $\cdot$ raw + regressors      & $\mathbf{2.7765}$ & $3.5824$ & \textbf{best}\\
D $\cdot$ raw, no regressors          & $3.0151$ & $3.7775$ & $+8.6\,\%$\\
C $\cdot$ compensated, no regressors  & $3.2906$ & $4.0972$ & $+18.5\,\%$\\
B $\cdot$ compensated + regressors    & $3.4718$ & $4.3762$ & $+25.0\,\%$\\
\bottomrule
\end{tabular}
\end{table}

\begin{figure}[H]
\includegraphics{S2_28_st02_experiment_neuralprophet.png}
\caption{One-hour-ahead forecast under NeuralProphet, 2021.}
\label{fig:np}
\end{figure}

Under the deployed model the ordering is the same as in the current era: both raw
configurations beat both compensated ones. The margins are smaller than the ridge
model's, as they should be --- NeuralProphet's own seasonal terms absorb part of what
the correction damages.

%% ==========================================================================
\section{Verdict}

\panel{Three criteria, three failures --- on independent data}{%
\textbf{1. Thermal dependence is not reduced.} $|r|$ with air temperature rises from
$0.238$ to $0.738$ and reverses sign.\par\medskip
\textbf{2. The coefficient is not near the variance-minimising value.} $5.00$ against
$1.00$ \mdegC, more than doubling residual variance relative to doing
nothing.\par\medskip
\textbf{3. The forecast does not improve.} The controlled-contrast penalty is
$+47.4\,\%$ at one hour and $+0.7\,\%$ at twenty-four, and $+18.5\,\%$ under
NeuralProphet. No horizon and no model favoured the compensated target.}

\subsection{Across both eras}

\begin{table}[H]
\centering
\caption{The same quantities, measured independently in each era.}
\begin{tabular}{@{}l R{2.6cm} R{2.6cm} R{2.4cm}@{}}
\toprule
\textbf{Quantity} & \textbf{Legacy 2021} & \textbf{Current era} & \textbf{Documented}\\
\midrule
Fitted \texttt{tair} slope, levels [\mdegC]      & $0.91$ & $1.63$ & $\mathbf{5.0}$\\
Fitted \texttt{tair} slope, differences [\mdegC] & $2.26$ & $2.17$ & $\mathbf{5.0}$\\
Variance-minimising coefficient [\mdegC]         & $1.00$ & $1.60$ & $\mathbf{5.0}$\\
Diurnal-band $R^2$, instantaneous \texttt{tair}  & $0.915$ & $0.932$ & ---\\
\bottomrule
\end{tabular}
\end{table}

Two things are worth noting in this table. The differenced slopes agree closely across
eras and instruments --- $2.26$ against $2.17$ \mdegC\ --- as do the diurnal-band $R^2$
values, $0.915$ against $0.932$. \textbf{The fast thermal response of this wall is a
stable, reproducible property.} The levels slopes do not agree, because they are
dominated by the seasonal behaviour, which differs between a half cycle and a full one.

Together with the three legacy stations measured from the supplied \texttt{Data.mat} ---
$7.21$, $1.33$ and $9.66$ \mdegC\ --- there are now five independent estimates of the
levels slope spanning $0.91$ to $9.66$. \textbf{No single constant serves them, and none
is near $5.0$.}

\subsection{Recommendation}

\textbf{Remove the fixed compensation from Notebook 00.} The recommendation is unchanged
from the current-era study and now rests on two independent datasets, two instruments,
two regressor sets and two model families.

Beyond that, this study says how the thermal response should be handled instead, and the
answer is that the two timescales must be separated:

\begin{itemize}
\item \textbf{The daily response} is contemporaneous, tight and reproducible
($R^2 = 0.92$ on the diurnal band, instantaneously). Instantaneous air temperature is an
excellent regressor for it, and no lag or time constant improves on that.
\item \textbf{The seasonal response} leads air temperature by roughly two months and
cannot be represented by any coefficient or lag on it. It belongs in the decomposition's
own annual term, or in a driver that carries the right phase. Installing a radiometer at
the legacy stations, or reconstructing irradiance on the wall's actual orientation from
a reanalysis product, is the measurement that would settle it.
\end{itemize}

%% ==========================================================================
\section{Caveats}

\begin{itemize}
\item One calendar year gives exactly one pass through the annual cycle, so the seasonal
lead is measured once and cannot be checked against a second cycle. Read it as a
description of 2021. Two estimators agree on its sign and order of magnitude but differ
on its size ($59$ against $89$ days), which is the honest precision available.
\item The attribution of the seasonal lead to solar geometry is an interpretation, not a
measurement. The legacy era has no radiometer; this study establishes only that air
temperature is the wrong driver at that timescale.
\item The trend and annual components of the decomposition are not separately
identifiable, as Section~4.1 shows explicitly. Neither is read as structural drift.
\item Only the regressors are imputed ($2.5\,\%$ of hours) and filled values are
systematically smoother than observed ones. The inclinometer keeps its gaps.
\item In the NeuralProphet experiment, gaps in the model input are bridged by
interpolation so the autoregressive window has continuous history; scoring happens only
at observed timestamps, identically for all configurations.
\item \texttt{lc\_lib.py} imports from \texttt{../thermal\_compensation/tc\_lib.py}.
Moving or renaming either study folder breaks that import.
\end{itemize}

%% ==========================================================================
\appendix
\section{Reproduction}

\begin{verbatim}
conda activate neuralprophet_env
cd studies/thermal_compensation_legacy
jupytext --to ipynb legacy_compensation_study.py
jupyter nbconvert --to notebook --execute --inplace \
        legacy_compensation_study.ipynb
\end{verbatim}

\noindent Edit the paired \texttt{.py} file, never the \texttt{.ipynb}. Setting
\texttt{CONTEXT = 'paper'} in the parameter cell re-exports every figure at manuscript
column size. The archive is never written to.

\medskip
\noindent\textbf{One library workaround is in force.} NeuralProphet 0.8.0 raises
\texttt{UnboundLocalError: conditional\_cols} when a prediction frame ends in
\texttt{NaN} and imputation is disabled --- the variable is bound only inside the
imputation branch but is read again when trailing rows are restored. The hourly grid
ends in a gap where the 20-minute grid did not, so \texttt{decompose\_series} trims the
trailing gap before predicting and reindexes the components back onto the full span,
leaving those hours empty.

\section{Output inventory}

\begin{table}[H]
\centering
\small
\caption{Artefacts produced by the notebook. Identifiers are stable names, not an
ordering.}
\begin{tabular}{@{}l L{9.2cm} c@{}}
\toprule
\textbf{Identifier} & \textbf{Contents} & \textbf{In this report}\\
\midrule
S2\_01 & Year completeness scan & Table 2\\
S2\_02 & Year completeness figure & Fig.~\ref{fig:years}\\
S2\_03 & Channel coverage, 2021 & Table 4\\
S2\_04 & Gap profile & Table 3\\
S2\_05 & Year overview & Fig.~\ref{fig:ov}\\
S2\_06 & Decomposition, three series & Figs.~\ref{fig:decinc}--\ref{fig:decrh}\\
S2\_07 & Decomposition summary & Table 5\\
S2\_08 & Imputation, \texttt{tair} and \texttt{rh} & Figs.~\ref{fig:imptair}, \ref{fig:imprh}\\
S2\_09, S2\_10 & Correlation matrices & Table 6\\
S2\_11 & Correlation heatmaps & Fig.~\ref{fig:heat}\\
S2\_12, S2\_31 & Cross-correlation peaks, differences and levels & Table 7\\
S2\_13 & Lag structure figure & Fig.~\ref{fig:ccf}\\
S2\_14, S2\_15 & Regression slopes & Table 10\\
S2\_16 & Fitted versus documented thermal slope & Fig.~\ref{fig:scatter}\\
S2\_17, S2\_18 & Acceptance tests and signed correlations & Tables 11, 12\\
S2\_19 & Effect of the documented compensation & Fig.~\ref{fig:effect}\\
S2\_20, S2\_21 & Coefficient sweep and figure & Table 13, Fig.~\ref{fig:sweep}\\
S2\_22, S2\_24 & Ridge results, 1 h and 24 h & Tables 14, 15\\
S2\_23, S2\_25 & Experiment figures & Figs.~\ref{fig:h1}, \ref{fig:h24}\\
S2\_26 & Controlled contrast & Table 16\\
S2\_27, S2\_28 & NeuralProphet results and figure & Table 17, Fig.~\ref{fig:np}\\
S2\_29 & Verdict & Section 9\\
S2\_30 & Cross-era comparison & Table 18\\
S2\_32--S2\_35 & Inertia scans and best time constants & Table 8\\
S2\_36, S2\_37 & Inertia scan figures & Figs.~\ref{fig:tau}, \ref{fig:tauband}\\
S2\_38 & Lag summary & Section~\ref{sec:lag}\\
S2\_39, S2\_41 & Seasonal lag and phase & Table 9\\
S2\_40 & Seasonal cross-correlation figure & Fig.~\ref{fig:ccfseason}\\
\bottomrule
\end{tabular}
\end{table}

\end{document}
